\documentclass[10pt]{article}

\usepackage[margin=0.72in]{geometry}
\usepackage[T1]{fontenc}
\usepackage[utf8]{inputenc}
\usepackage{lmodern}
\usepackage{enumitem}
\usepackage{titlesec}
\usepackage{xcolor}
\usepackage{hyperref}

\hypersetup{
  colorlinks=true,
  linkcolor=blue!50!black,
  urlcolor=blue!50!black
}

\setlist[itemize]{leftmargin=*, itemsep=2pt, topsep=3pt}
\setlist[enumerate]{leftmargin=*, itemsep=2pt, topsep=3pt}
\titlespacing*{\section}{0pt}{8pt}{4pt}
\titlespacing*{\subsection}{0pt}{6pt}{3pt}
\titleformat{\section}{\large\bfseries}{\thesection}{0.5em}{}
\titleformat{\subsection}{\normalsize\bfseries}{\thesubsection}{0.5em}{}
\pagestyle{plain}

\title{\vspace{-2.0em}\textbf{SPR 4862 High-Level Implementation Outline}\\
\large INDOT Solar Development Suitability Map}
\author{Md Kamruzzaman Kamrul}
\date{May 19, 2026}

\begin{document}
\maketitle
\vspace{-1.2em}

\section{Project Goal}

The goal of this project is to develop a public-facing, interactive web map for
INDOT solar development suitability. The map will use the provided GIS shapefiles
for candidate sites, INDOT facilities, and right-of-way/interchange sites. The
system will allow users to inspect site locations, view relevant suitability
attributes, filter sites, and support future data updates through a repeatable
processing pipeline.

\section{Proposed Implementation Approach}

\subsection{1. Data Preparation and Validation}
\begin{itemize}
  \item Inspect the provided shapefiles and document the available fields, record counts, coordinate systems, and geometry types.
  \item Clean the spatial data before it is used in the web application.
  \item Convert the GIS data into web-friendly formats such as GeoJSON and CSV.
  \item Preserve original source fields while adding cleaned display fields for the public map.
\end{itemize}

\subsection{2. Data Processing Pipeline}
\begin{itemize}
  \item Build a repeatable Python-based pipeline to read, validate, clean, and export the site data.
  \item Repair geometry issues where needed.
  \item Derive usable latitude and longitude from the actual geometries.
  \item Normalize fields so that values are displayed consistently across layers.
\end{itemize}

\subsection{3. Web Map Prototype}
\begin{itemize}
  \item Develop an initial interactive map using cleaned GeoJSON data.
  \item Add separate layers for facility and ROW/interchange sites.
  \item Add basic popups, legends, and site-type filters.
  \item Use this prototype to confirm the data model and visual design before building a larger backend.
\end{itemize}

\subsection{4. Backend and Configuration}
\begin{itemize}
  \item After the cleaned data model is stable, add a backend API for site queries, details, statistics, and exports.
  \item Use a configuration file to control display labels, visible fields, filters, and layer behavior.
  \item Keep the system maintainable so future shapefile updates can be processed with minimal manual work.
\end{itemize}

\subsection{5. Final Application, Testing, and Documentation}
\begin{itemize}
  \item Integrate the frontend, backend, and cleaned data outputs.
  \item Test map performance, browser behavior, filtering, popups, and data exports.
  \item Prepare setup documentation, data-update instructions, and a final technical summary.
\end{itemize}

\section{Initial Assessment Completed}

I completed an initial inspection of the local GIS files and generated a data
inventory report. The initial assessment confirms that the main shapefiles are
available and readable:

\begin{itemize}
  \item Candidate sites: 104 records.
  \item Facility scored sites: 58 records.
  \item ROW/interchange scored sites: 45 records.
  \item All three datasets use the same projected coordinate system, EPSG:26916.
  \item The data can be read successfully in Python using GeoPandas.
\end{itemize}

This means the project can move forward with a Python-based data pipeline and a
web mapping workflow. The next immediate technical step is to create the cleaned
export files that the first web map prototype will consume.

\section{Current Problems I Am Working Through}

During the initial assessment, I found several data issues that need to be handled
before the map is built. I am currently working on solving these as part of the
data-cleaning pipeline.

\begin{itemize}
  \item The latitude and longitude fields in the shapefiles are not usable because they are stored as zero values. I will derive real map coordinates from the geometry instead.
  \item One candidate site appears in the all-candidate file but does not appear in the scored facility or ROW files. I am checking how to represent this site without losing it.
  \item One geometry has a self-intersection issue. I will repair this during export so it does not cause problems in GeoJSON or PostGIS.
  \item Some land-cover fields are stored as percentages in one file and fractions in the scored files. I will normalize these values for consistent display.
  \item Some score fields appear to be layer-specific, meaning they should not be shown as equally applicable for every site type.
  \item Some site labels and site-type values need cleanup for public display while preserving the original source data.
\end{itemize}

\section{Expected Near-Term Output}

The next working output will be a cleaned dataset package containing:

\begin{itemize}
  \item Clean GeoJSON files for the first map prototype.
  \item A CSV export for easier review.
  \item A short validation report showing row counts, repaired geometries, missing score records, and normalized fields.
  \item A simple interactive map prototype to verify that the cleaned data displays correctly.
\end{itemize}

\section{Requested Direction and Feedback}

At this stage, I am working based on my current experience, available materials,
and technical understanding of the GIS data and web-mapping workflow. I will share
the working materials and updates momentarily as I continue to organize the project.
If any part of the plan, data handling approach, or implementation direction should
be changed, please let me know so I can address it early and adjust the work.

If you have any specific directives for the map, interface, data fields, or expected
workflow, I can try to incorporate those as well. I will also maintain the project
through GitHub and update it from time to time as the implementation progresses.
GitHub repository link: \url{https://github.com/kamrul28890/InDoT-map-project}.

You also mentioned that Illinois had a similar interactive map, but the original
public link does not appear to be available anymore. If there is any way to find,
access, or recover that map or related screenshots/documentation, could you please
look into it? Having access to the Illinois example would help significantly in
understanding the expected requirements, interface behavior, and stakeholder-facing
features for the INDOT version.

\end{document}
